\section{Related Work}
\label{sec:related}

\para{Internal-state probing for truthfulness.}
A growing body of work trains classifiers on LLM hidden-layer activations to predict statement truthfulness. \citet{azaria2023internal} propose SAPLMA, which achieves 71--83\% accuracy using feed-forward classifiers on hidden states, demonstrating that models encode truth-relevant information internally. \citet{burns2022discovering} introduce Contrast-Consistent Search (CCS), an unsupervised method that extracts latent knowledge from contrast pairs without labels, achieving robustness to misleading prompts. \citet{marks2023geometry} show that true and false statements occupy geometrically separated regions in representation space, forming linear ``truth directions.'' \citet{li2023representation} develop Representation Engineering (RepE) for reading and controlling high-level cognitive properties including truthfulness. These methods detect \emph{whether} a statement is false but do not categorize \emph{why}---they treat all falsehoods as a single class. Unlike these approaches, we classify false outputs by their underlying mechanism without requiring access to internal representations.

\para{Hallucination detection and taxonomy.}
\citet{farquhar2024detecting} introduce semantic entropy, which clusters sampled responses by meaning equivalence and computes entropy over clusters, specifically targeting confabulations. This method is principled for epistemic uncertainty but by design cannot detect high-confidence errors where the model consistently produces the wrong answer. \citet{simhi2025hack} propose the HACK framework, which classifies hallucinations along Knowledge and Certainty axes, identifying Certainty Misalignment (CM) cases where models ``know'' the correct answer but confidently produce wrong ones. They validate this distinction using activation steering, showing differential effectiveness on knowledge-present versus knowledge-absent errors. HACK is the closest prior work to ours, but it does not examine incentive-driven misreporting (sycophancy) as a separate mechanism. \citet{sui2024confabulation} show that confabulated outputs display increased narrativity and semantic coherence, providing evidence that different falsehood types have measurable stylistic signatures.

\para{Sycophancy and incentive-driven behavior.}
\citet{sharma2024towards} systematically demonstrate that sycophancy is universal across production LLMs and trace its cause to the RLHF training pipeline, where human preference data rewards agreement with users. They show that user beliefs can shift model accuracy by up to 27 percentage points, establishing sycophancy as a canonical example of incentive-driven misreporting. \citet{hubinger2024sleeper} demonstrate that LLMs can be trained to exhibit deceptive behavior that persists through safety training, motivating the need for mechanism-specific detectors. \citet{hagendorff2023deception} show that LLMs can understand and induce deception strategies, establishing the capability for strategic deception beyond passive hallucination. This thread of research measures behavioral outcomes but does not use internal-state methods or quantify how sycophancy-driven errors relate to epistemic failures in prevalence.

\para{Benchmarks for LLM truthfulness.}
\citet{lin2021truthfulqa} introduce \truthfulqa, a benchmark of 817 questions designed to elicit false answers learned from imitating human misconceptions. While widely used, TruthfulQA does not distinguish between falsehood mechanisms: a model answering incorrectly could be repeating training data, genuinely uncertain, or exhibiting some other failure mode. \citet{chen2023llm} show that LLM-generated misinformation is harder to detect than human-written misinformation with the same semantics, and \citet{duan2024llms} demonstrate that LLMs react differently in hidden states when processing genuine versus fabricated responses. These works motivate the need for benchmarks that label errors by mechanism, which our framework provides.
